\section*{Разбор задачи из тестового варианта №3}
\addcontentsline{toc}{section}{Разбор задачи из тестового варианта №3}
\begin{Task}{Гауссовские векторы и условное ожидание нелинейной функции}{}
    Случайный вектор $(X, Y, Z)$ является гауссовским с нулевым средним и матрицей ковариаций
    \[
    \Sigma = \begin{pmatrix} 
    3 & -1 & 1 \\ 
    -1 & 13/3 & -1/3 \\ 
    1 & -1/3 & 10/3 
    \end{pmatrix}.
    \]
    Найдите $\E(\cos(YZ)|X)$.
\end{Task}

\begin{Solution}
    \textbf{Семантический анализ задачи.} Прямое интегрирование нелинейной функции $\cos(YZ)$ по двумерной условной плотности ведет к вычислительному коллапсу. Фундаментальный подход заключается в декомпозиции гауссовского вектора методом ортогональных проекций с последующим переходом в комплексную плоскость для вычисления математического ожидания тригонометрической функции.

    \begin{MethodBox}[Шаг 1. Ортогональное проецирование]
        Согласно теореме о нормальной корреляции, компоненты $Y$ и $Z$ можно представить в виде суммы линейной проекции на условие $X$ и независимого гауссовского остатка:
        \[
        Y = aX + \widetilde{Y}, \quad Z = bX + \widetilde{Z}, \quad \text{где } \widetilde{Y} \indep X, \widetilde{Z} \indep X.
        \]
    \end{MethodBox}

    Извлечем необходимые элементы из матрицы ковариаций $\Sigma$:
    \begin{itemize}
        \item Дисперсия условия: $\Var(X) = \Sigma_{11} = 3$.
        \item Ковариации с условием: $\Cov(Y, X) = \Sigma_{21} = -1$ и $\Cov(Z, X) = \Sigma_{31} = 1$.
        \item Внутренние характеристики: $\Var(Y) = \frac{13}{3}$, $\Var(Z) = \frac{10}{3}$, $\Cov(Y, Z) = -\frac{1}{3}$.
    \end{itemize}

    Вычислим коэффициенты проекций:
    \begin{align*}
        a &= \frac{\Cov(Y, X)}{\Var(X)} = -\frac{1}{3}, \\
        b &= \frac{\Cov(Z, X)}{\Var(X)} = \frac{1}{3}.
    \end{align*}

    Исследуем свойства ортогональных остатков $\widetilde{Y}$ и $\widetilde{Z}$. Поскольку исходный вектор имеет нулевое среднее, $\E(\widetilde{Y}) = \E(\widetilde{Z}) = 0$. Вычислим их ковариацию:
    \begin{align*}
        \Cov(\widetilde{Y}, \widetilde{Z}) &= \Cov(Y - aX, Z - bX) = \Cov(Y, Z) - ab\Var(X) \\
        &= -\frac{1}{3} - \left(-\frac{1}{3}\right)\left(\frac{1}{3}\right) \cdot 3 = -\frac{1}{3} + \frac{1}{3} = 0.
    \end{align*}
    В силу нормальности распределения некоррелированность означает независимость, следовательно, $\widetilde{Y} \indep \widetilde{Z}$. 
    
    Найдем дисперсии остатков:
    \begin{align*}
        \Var(\widetilde{Y}) &= \Var(Y) - a^2\Var(X) = \frac{13}{3} - \left(-\frac{1}{3}\right)^2 \cdot 3 = \frac{13}{3} - \frac{1}{3} = 4, \\
        \Var(\widetilde{Z}) &= \Var(Z) - b^2\Var(X) = \frac{10}{3} - \left(\frac{1}{3}\right)^2 \cdot 3 = \frac{10}{3} - \frac{1}{3} = 3.
    \end{align*}

    Таким образом, при фиксированном (условном) значении $X = x$, случайные величины $Y$ и $Z$ распадаются на две \textbf{независимые} гауссовские величины:
    \[
    Y_x \sim \N\left(-\frac{x}{3}, 4\right), \quad Z_x \sim \N\left(\frac{x}{3}, 3\right).
    \]

    \begin{HackBox}[Шаг 2. Комплексификация и итерированное ожидание]
        Интеграл от $\cos(Y_x Z_x)$ не вычисляется в действительной области. Перейдем к формуле Эйлера: $\cos(\varphi) = \text{Re}(e^{i\varphi})$. Применим свойство башни (итерированного математического ожидания), обусловив выражение относительно одной из независимых величин.
    \end{HackBox}

    Запишем искомое ожидание в комплексной форме:
    \[
    \E[\cos(Y_x Z_x)] = \text{Re}\left( \E[e^{i Y_x Z_x}] \right).
    \]
    Зафиксируем $Z_x$. Внутреннее ожидание представляет собой характеристическую функцию нормальной величины $Y_x$ в точке $t = Z_x$. Известно, что для $W \sim \N(\mu, \sigma^2)$ характеристическая функция равна $\E[e^{itW}] = \exp(i\mu t - \frac{1}{2}\sigma^2 t^2)$.
    \[
    \E_{Y_x}[e^{i Z_x Y_x} | Z_x] = \exp\left( i \left(-\frac{x}{3}\right) Z_x - \frac{1}{2}(4) Z_x^2 \right) = \exp\left( -i\frac{x}{3} Z_x - 2 Z_x^2 \right).
    \]

    Теперь необходимо вычислить внешнее ожидание по $Z_x \sim \N\left(\frac{x}{3}, 3\right)$. Задача сводится к вычислению интеграла Гаусса вида $\E[\exp(c Z_x + d Z_x^2)]$, где $c = -i\frac{x}{3}$ и $d = -2$.

    Для случайной величины $\eta \sim \N(\mu, \sigma^2)$ и произвольных констант $c, d$ (при $\text{Re}(1-2d\sigma^2)>0$) справедливо тождество:
    \[
    \E[\exp(c\eta + d\eta^2)] = \frac{1}{\sqrt{1 - 2d\sigma^2}} \exp\left( \frac{c^2\sigma^2 + 2c\mu + 2d\mu^2}{2(1 - 2d\sigma^2)} \right).
    \]

    Подставим параметры распределения $Z_x$ ($\mu = \frac{x}{3}$, $\sigma^2 = 3$):
    \begin{itemize}
        \item Знаменатель под корнем: $1 - 2d\sigma^2 = 1 - 2(-2)(3) = 1 + 12 = 13$.
        \item Элементы числителя экспоненты:
        \begin{align*}
            c^2\sigma^2 &= \left(-i\frac{x}{3}\right)^2 \cdot 3 = \left(-\frac{x^2}{9}\right) \cdot 3 = -\frac{3x^2}{9}, \\
            2c\mu &= 2\left(-i\frac{x}{3}\right)\left(\frac{x}{3}\right) = -i\frac{2x^2}{9}, \\
            2d\mu^2 &= 2(-2)\left(\frac{x}{3}\right)^2 = -\frac{4x^2}{9}.
        \end{align*}
    \end{itemize}

    Суммируем числитель:
    \[
    c^2\sigma^2 + 2c\mu + 2d\mu^2 = \frac{-3x^2 - 4x^2 - i 2x^2}{9} = \frac{-7x^2 - i 2x^2}{9}.
    \]

    Формируем итоговый аргумент экспоненты, разделив на $2(1 - 2d\sigma^2) = 26$:
    \[
    \frac{1}{26} \left( \frac{-7x^2 - i 2x^2}{9} \right) = -\frac{7x^2}{234} - i\frac{2x^2}{234} = -\frac{7}{234}x^2 - i\frac{1}{117}x^2.
    \]

    Собираем комплексное ожидание:
    \[
    \E[e^{i Y_x Z_x}] = \frac{1}{\sqrt{13}} \exp\left( -\frac{7}{234}x^2 \right) \exp\left( -i\frac{1}{117}x^2 \right).
    \]

    Для получения окончательного результата выделяем действительную часть выражения, используя формулу Эйлера $\exp(-i\alpha) = \cos(\alpha) - i\sin(\alpha)$:
    \[
    \text{Re}\left( \E[e^{i Y_x Z_x}] \right) = \frac{1}{\sqrt{13}} \exp\left( -\frac{7}{234}x^2 \right) \cos\left( \frac{1}{117}x^2 \right).
    \]

    Остается произвести обратную замену детерминированного $x$ на случайную величину $X$ в силу свойств условного математического ожидания.

\end{Solution}

\begin{Answer}
    \[
    \E(\cos(YZ)|X) = \frac{1}{\sqrt{13}} \exp\left( -\frac{7}{234}X^2 \right) \cos\left( \frac{1}{117}X^2 \right).
    \]
\end{Answer}
